\documentclass[12pt,a4paper,twoside]{report}

% ─── Encoding & Language ───
\usepackage[utf8]{inputenc}
\usepackage[T1]{fontenc}
\usepackage[english]{babel}

% ─── Page Geometry ───
\usepackage[top=2.5cm, bottom=2.5cm, left=3cm, right=2.5cm]{geometry}

% ─── Fonts ───
\usepackage{lmodern}
\usepackage{microtype}

% ─── Math ───
\usepackage{amsmath,amssymb,amsfonts,mathtools}
\usepackage{bm}

% ─── Graphics & Figures ───
\usepackage[dvipsnames,svgnames,x11names]{xcolor}
\usepackage{graphicx}
\usepackage{svg}
\usepackage{float}
\usepackage{subcaption}

% ─── Tables ───
\usepackage{booktabs}
\usepackage{array}
\usepackage{multirow}
\usepackage{tabularx}

% ─── Code Listings ───
\usepackage{listings}
\lstset{
  basicstyle=\ttfamily\small,
  keywordstyle=\color{blue}\bfseries,
  commentstyle=\color{gray},
  stringstyle=\color{red},
  breaklines=true,
  frame=single,
  numbers=left,
  numberstyle=\tiny\color{gray},
}

% ─── Algorithms ───
\usepackage{algorithm}
\usepackage{algorithmic}

% ─── Hyperlinks ───
\usepackage[colorlinks=true, linkcolor=NavyBlue, citecolor=OliveGreen, urlcolor=BrickRed]{hyperref}

% ─── Bibliography ───
\usepackage[numbers,sort&compress]{natbib}

% ─── Headers & Footers ───
\usepackage{fancyhdr}
\pagestyle{fancy}
\fancyhf{}
\fancyhead[LE]{\leftmark}
\fancyhead[RO]{\rightmark}
\fancyfoot[C]{\thepage}
\renewcommand{\headrulewidth}{0.4pt}

% ─── TikZ for diagrams ───
\usepackage{tikz}
\usetikzlibrary{arrows.meta, positioning, shapes.geometric, calc}

% ─── Theorem environments ───
\usepackage{amsthm}
\newtheorem{definition}{Definition}[chapter]
\newtheorem{remark}{Remark}[chapter]

% ─── Custom commands ───
\newcommand{\vect}[1]{\bm{#1}}
\newcommand{\mat}[1]{\mathbf{#1}}
\newcommand{\R}{\mathbb{R}}
\newcommand{\dd}{\mathrm{d}}
\newcommand{\pd}[2]{\frac{\partial #1}{\partial #2}}
\newcommand{\norm}[1]{\left\| #1 \right\|}

% ─── Spacing ───
\usepackage{setspace}
\onehalfspacing
\usepackage{parskip}

% ─── Glossary / Acronyms ───
\usepackage[acronym,toc]{glossaries}
\makeglossaries
\newacronym{iga}{IGA}{Isogeometric Analysis}
\newacronym{rom}{ROM}{Reduced Order Model}
\newacronym{fom}{FOM}{Full Order Model}
\newacronym{pod}{POD}{Proper Orthogonal Decomposition}
\newacronym{svd}{SVD}{Singular Value Decomposition}
\newacronym{nurbs}{NURBS}{Non-Uniform Rational B-Splines}
\newacronym{fem}{FEM}{Finite Element Method}
\newacronym{fsi}{FSI}{Fluid--Structure Interaction}
\newacronym{lhs}{LHS}{Latin Hypercube Sampling}
\newacronym{deim}{DEIM}{Discrete Empirical Interpolation Method}
\newacronym{ecsw}{ECSW}{Energy-Conserving Sampling and Weighting}
\newacronym{dt}{DT}{Digital Twin}
\newacronym{cnam}{Cnam}{Conservatoire National des Arts et Métiers}

\begin{document}

% ════════════════════════════════════════════════════════════════
%  FRONT MATTER
% ════════════════════════════════════════════════════════════════
% ════════════════════════════════════════════════════════════════
%  TITLE PAGE
% ════════════════════════════════════════════════════════════════
\begin{titlepage}
\centering

\vspace*{1cm}

% ─── Institution logos ───
\includesvg[height=2.5cm]{images/CNAM_Logo}
% \hspace{2cm}
% \includegraphics[height=2.5cm]{figures/epn04_logo.png}

\vspace{1cm}

{\Large\scshape Conservatoire National des Arts et Métiers}\\[0.3cm]
{\large Équipement pour la Performance et le Numérique (EPN04)}

\vspace{1.5cm}

\rule{\textwidth}{1.5pt}\\[0.6cm]
{\Huge\bfseries Digital Twin \& Reduced Order Modelling\\[0.4cm] for Nonlinear Dynamics}\\[0.4cm]
\rule{\textwidth}{1.5pt}

\vspace{1.5cm}

{\Large\itshape Internship Report}\\[0.5cm]
{\large March -- September 2026}

\vspace{2cm}

\begin{minipage}[t]{0.45\textwidth}
\begin{flushleft}
\textbf{Author:}\\
Mumen Wehbe
\end{flushleft}
\end{minipage}
\hfill
\begin{minipage}[t]{0.45\textwidth}
\begin{flushright}
\textbf{Supervisor:}\\
Dr.\ Christophe Hoareau\\
EPN04, Cnam
\end{flushright}
\end{minipage}

\vfill

{\large \the\year}

\end{titlepage}


\pagenumbering{roman}

% Abstract
\chapter*{Abstract}
\addcontentsline{toc}{chapter}{Abstract}
This report presents the work conducted during a six-month internship at the \acrfull{cnam}, within the Équipement pour la Performance et le Numérique (EPN04) department, under the supervision of Dr.\ Christophe Hoareau. The project addresses the development of a real-time pedagogical Digital Twin for nonlinear structural mechanics, combining \acrfull{iga} with projection-based \acrfull{rom} techniques.

The methodology relies on \acrfull{nurbs} basis functions for exact geometric representation and employs \acrfull{pod} via \acrfull{svd} to extract dominant deformation modes from parametric snapshot databases. A Galerkin projection framework reduces the dimensionality of the governing equations, while hyper-reduction strategies (DEIM, ECSW) accelerate the evaluation of nonlinear terms.

The resulting web-based application enables real-time simulation of geometrically nonlinear 2D structures with parametric control, achieving solve times under 10\,ms compared to several seconds for the full-order model. The platform serves as both a research validation tool and an interactive pedagogical environment for teaching advanced computational mechanics.

\vspace{1cm}
\noindent\textbf{Keywords:} Isogeometric Analysis, Reduced Order Modelling, Nonlinear Mechanics, Digital Twin, NURBS, Proper Orthogonal Decomposition, Hyper-reduction, Real-time Simulation.

% Table of Contents
\tableofcontents
\listoffigures
\listoftables
\printglossary[type=\acronymtype, title=List of Acronyms]

\cleardoublepage
\pagenumbering{arabic}

% ════════════════════════════════════════════════════════════════
%  MAIN MATTER
% ════════════════════════════════════════════════════════════════
% ════════════════════════════════════════════════════════════════
%  INTRODUCTION
% ════════════════════════════════════════════════════════════════
\chapter*{Introduction}
\addcontentsline{toc}{chapter}{Introduction}
\markboth{Introduction}{Introduction}

% ─────────────────────────────────────────────────────────────
\section*{Context}
% ─────────────────────────────────────────────────────────────

The rapid advancement of computational mechanics and the growing demand for interactive simulation tools have created a fertile ground for innovation at the intersection of numerical methods and digital engineering. This internship takes place within the Équipement pour la Performance et le Numérique (EPN04) department at the \acrlong{cnam} in Paris, under the supervision of Dr.\ Christophe Hoareau.

\paragraph{Who.} The project is carried out by WEHBE Mumen, an engineering student, in collaboration with the EPN04 research team specializing in advanced numerical simulation and structural mechanics.

\paragraph{What.} The objective is to develop a real-time \acrfull{dt} platform that couples a physical 3D-printed structural model with a virtual simulation environment powered by \acrfull{rom} techniques built upon \acrfull{iga}.

\paragraph{Where.} The work is conducted at the Cnam campus (Paris), leveraging the department's computational infrastructure and pedagogical environment.

\paragraph{When.} The internship spans six months, from March to September 2026, following a structured timeline that progresses from theoretical foundations through computational implementation to application deployment.

\paragraph{Why.} Traditional high-fidelity simulations of nonlinear structural problems are computationally expensive, often requiring minutes to hours for a single parameter evaluation. This renders them impractical for real-time applications, interactive teaching, and rapid design exploration. The need for a methodology that preserves simulation fidelity while achieving real-time performance motivates this work.

% ─────────────────────────────────────────────────────────────
\section*{Problematic and Research Gap}
% ─────────────────────────────────────────────────────────────

While \acrlong{rom} techniques have been extensively studied for linear structural problems, their application to \textbf{geometrically nonlinear} problems with \textbf{parametric variations} remains an active area of research. The key challenges include:

\begin{itemize}
    \item \textbf{Nonlinear operator dependence:} The tangent stiffness matrix $\mat{K}_T(\vect{u})$ changes at every Newton--Raphson iteration, requiring repeated projection onto the reduced basis---a computational bottleneck that standard Galerkin projection alone cannot resolve efficiently.
    
    \item \textbf{Geometric parameterization:} When the geometry itself is a parameter (e.g., varying aspect ratios or hole sizes), the physical domain $\Omega(\mu)$ changes, invalidating the fixed-mesh assumption of classical ROM approaches.
    
    \item \textbf{Hyper-reduction for nonlinear terms:} Even after projection, evaluating internal forces $\mat{F}_{\text{int}}(\vect{u})$ requires integration over the full mesh. Hyper-reduction techniques (DEIM, ECSW) are needed to reduce this cost, but their application to IGA-discretized nonlinear problems is not yet well established.
    
    \item \textbf{Pedagogical accessibility:} Existing ROM implementations are typically confined to research codes (Python/MATLAB) that are inaccessible to students without specialized software environments.
\end{itemize}

% ─────────────────────────────────────────────────────────────
\section*{Objectives}
% ─────────────────────────────────────────────────────────────

The primary objectives of this internship are:

\begin{enumerate}
    \item \textbf{Develop an IGA-based computational core} capable of solving 2D geometrically nonlinear elasticity problems with NURBS discretization, including stiffness assembly, Newton--Raphson iteration, and stress recovery.
    
    \item \textbf{Construct a projection-based ROM framework} using POD/SVD for basis extraction and Galerkin projection for operator reduction, extended to handle nonlinear terms via hyper-reduction.
    
    \item \textbf{Build a real-time web application} that serves as both a research validation tool and an interactive pedagogical platform, enabling students to explore IGA, ROM, and nonlinear mechanics concepts through live simulations.
    
    \item \textbf{Validate the methodology} through benchmark problems (cantilever beam, plate with hole) and quantify the speedup and accuracy of the ROM against the full-order model.
    
    \item \textbf{Explore 3D extensions} (time permitting) to assess the scalability of the approach to industrially relevant three-dimensional geometries.
\end{enumerate}

% ════════════════════════════════════════════════════════════════
%  CHAPTER 1 : STATE OF THE ART
% ════════════════════════════════════════════════════════════════
\chapter{State of the Art}
\label{ch:state-of-art}

This chapter reviews the theoretical foundations underpinning the three pillars of the project: Isogeometric Analysis, continuum mechanics for linear and nonlinear elasticity, and projection-based reduced order modelling.

% ──────────────────────────────────────────────────────────────
\section{Isogeometric Analysis (IGA)}
\label{sec:iga}
% ──────────────────────────────────────────────────────────────

\subsection{From FEM to IGA}

The \acrfull{fem} has been the dominant discretization method in computational mechanics for over half a century. However, a fundamental inefficiency persists in the traditional FEM workflow: the geometry created in Computer-Aided Design (CAD) systems using \acrfull{nurbs} must be re-approximated by polynomial meshes before analysis can begin. This mesh generation step introduces geometric errors, is labor-intensive, and often constitutes a significant portion of the total simulation time.

\acrfull{iga}, introduced by Hughes, Cottrell, and Bazilevs in 2005~\cite{hughes2005}, resolves this gap by using the same NURBS basis functions for both geometry representation and solution approximation. The key principle is:
\begin{equation}
    \vect{x}(\xi) = \sum_{i=1}^{n} R_{i,p}(\xi) \, \vect{P}_i
    \label{eq:iga-mapping}
\end{equation}
where $R_{i,p}$ are the rational basis functions and $\vect{P}_i$ are the control points. The displacement field is approximated using the same basis:
\begin{equation}
    \vect{u}^h(\xi) = \sum_{i=1}^{n} R_{i,p}(\xi) \, \vect{d}_i
\end{equation}
where $\vect{d}_i$ are the unknown displacement coefficients at control points.

\subsection{B-Splines and NURBS}

A B-spline of degree $p$ is defined over a knot vector $\Xi = \{\xi_1, \xi_2, \ldots, \xi_{n+p+1}\}$ by the Cox--de Boor recursion:
\begin{align}
    N_{i,0}(\xi) &= \begin{cases} 1 & \text{if } \xi_i \leq \xi < \xi_{i+1} \\ 0 & \text{otherwise} \end{cases} \\[6pt]
    N_{i,p}(\xi) &= \frac{\xi - \xi_i}{\xi_{i+p} - \xi_i} N_{i,p-1}(\xi) + \frac{\xi_{i+p+1} - \xi}{\xi_{i+p+1} - \xi_{i+1}} N_{i+1,p-1}(\xi)
\end{align}

NURBS basis functions are obtained by introducing weights $w_i > 0$:
\begin{equation}
    R_{i,p}(\xi) = \frac{N_{i,p}(\xi) \, w_i}{\sum_{j=1}^{n} N_{j,p}(\xi) \, w_j}
    \label{eq:nurbs-basis}
\end{equation}

For 2D problems, tensor-product surfaces are constructed:
\begin{equation}
    \vect{S}(\xi, \eta) = \sum_{i=1}^{n} \sum_{j=1}^{m} R_{i,j}^{p,q}(\xi, \eta) \, \vect{P}_{i,j}
\end{equation}

\subsection{Refinement Strategies}

IGA supports three refinement mechanisms that preserve the exact geometry:
\begin{itemize}
    \item \textbf{$h$-refinement (knot insertion):} Adds new knots to the knot vector, increasing the number of basis functions without changing the polynomial degree.
    \item \textbf{$p$-refinement (order elevation):} Increases the polynomial degree while maintaining the knot vector structure.
    \item \textbf{$k$-refinement:} A combination unique to IGA---order elevation followed by knot insertion---that produces basis functions with maximal continuity ($C^{p-1}$).
\end{itemize}

\subsection{Advantages of IGA for ROM}

Several properties make IGA particularly well-suited for reduced order modelling:
\begin{enumerate}
    \item \textbf{Exact geometry:} No geometric approximation error, critical for parametric studies where the domain changes.
    \item \textbf{Higher continuity:} $C^{p-1}$ inter-element continuity (vs.\ $C^0$ in FEM) produces smoother solution fields, leading to better snapshot quality.
    \item \textbf{Natural parameterization:} The mapping from the master domain $\hat{\Omega}$ to the physical domain $\Omega(\mu)$ is built into the NURBS framework, enabling parametric ROM on a fixed reference configuration.
\end{enumerate}

% ──────────────────────────────────────────────────────────────
\section{Linear and Nonlinear Elasticity}
\label{sec:elasticity}
% ──────────────────────────────────────────────────────────────

\subsection{Linear Elasticity}

Under the assumption of infinitesimal strains, the linearized strain tensor is:
\begin{equation}
    \bm{\varepsilon} = \frac{1}{2}\left(\nabla\vect{u} + \nabla\vect{u}^T\right)
\end{equation}

The constitutive law for an isotropic material (generalized Hooke's law) in Voigt notation reads:
\begin{equation}
    \bm{\sigma} = \mat{C} \, \bm{\varepsilon}
\end{equation}
where $\mat{C}$ is the elasticity matrix depending on Young's modulus $E$ and Poisson's ratio $\nu$.

The weak form of the equilibrium equation leads to the standard stiffness system:
\begin{equation}
    \mat{K} \, \vect{u} = \vect{F}
    \label{eq:linear-system}
\end{equation}

\subsection{Geometric Nonlinearity}

When displacements become large relative to the structural dimensions, the linear strain measure is no longer adequate. The deformation is described by the deformation gradient:
\begin{equation}
    \mat{F} = \pd{\vect{x}}{\vect{X}} = \mat{I} + \nabla_{\vect{X}} \vect{u}
    \label{eq:deformation-gradient}
\end{equation}

The \textbf{Green--Lagrange strain tensor} captures finite deformations without contamination by rigid-body rotations:
\begin{equation}
    \mat{E} = \frac{1}{2}\left(\mat{F}^T \mat{F} - \mat{I}\right) = \frac{1}{2}\left(\nabla\vect{u} + \nabla\vect{u}^T + \nabla\vect{u}^T \nabla\vect{u}\right)
    \label{eq:green-lagrange}
\end{equation}

The quadratic term $\nabla\vect{u}^T \nabla\vect{u}$ is the source of geometric nonlinearity, producing the \textbf{geometric stiffening} effect.

\subsection{Constitutive Model: Saint Venant--Kirchhoff}

For this work, we adopt the Saint Venant--Kirchhoff (SVK) hyperelastic model, which relates the second Piola--Kirchhoff stress tensor $\mat{S}$ to the Green--Lagrange strain:
\begin{equation}
    \mat{S} = \mat{C} : \mat{E}
\end{equation}

This is the simplest hyperelastic model and is appropriate for large displacements with moderate strains.

\subsection{Newton--Raphson Solution Procedure}

The nonlinear equilibrium is expressed as a residual equation:
\begin{equation}
    \vect{R}(\vect{u}) = \vect{F}_{\text{ext}} - \vect{F}_{\text{int}}(\vect{u}) = \vect{0}
\end{equation}

Linearization yields the iterative correction:
\begin{equation}
    \mat{K}_T(\vect{u}^{(k)}) \, \Delta\vect{u} = \vect{R}(\vect{u}^{(k)}), \qquad \vect{u}^{(k+1)} = \vect{u}^{(k)} + \Delta\vect{u}
\end{equation}
where the tangent stiffness $\mat{K}_T = \mat{K}_L + \mat{K}_G$ includes both the material (linear) and geometric contributions.

% ──────────────────────────────────────────────────────────────
\section{Projection-Based Reduced Order Modelling}
\label{sec:rom}
% ──────────────────────────────────────────────────────────────

\subsection{The Offline--Online Paradigm}

ROM follows a two-phase strategy:
\begin{itemize}
    \item \textbf{Offline phase (expensive, done once):} Generate high-fidelity solutions (snapshots) for a training set of parameters, extract a reduced basis, and pre-compute projected operators.
    \item \textbf{Online phase (fast, repeated):} For any new parameter value, solve the reduced system of size $r \ll N$ in real time.
\end{itemize}

\subsection{Proper Orthogonal Decomposition (POD)}

Given a snapshot matrix $\mat{S} = [\vect{u}(\mu_1), \ldots, \vect{u}(\mu_{N_s})] \in \R^{N \times N_s}$, the POD basis is obtained via the Singular Value Decomposition:
\begin{equation}
    \mat{S} = \mat{U} \, \mat{\Sigma} \, \mat{V}^T
\end{equation}

The reduced basis $\mat{\Phi} \in \R^{N \times r}$ consists of the first $r$ columns of $\mat{U}$, chosen such that the cumulative energy criterion is satisfied:
\begin{equation}
    \eta(r) = \frac{\sum_{i=1}^{r} \sigma_i^2}{\sum_{i=1}^{N_s} \sigma_i^2} > 1 - \varepsilon, \qquad \text{typically } \varepsilon = 10^{-4}
    \label{eq:energy-criterion}
\end{equation}

\subsection{Galerkin Projection}

The displacement is approximated in the reduced subspace:
\begin{equation}
    \vect{u} \approx \mat{\Phi} \, \vect{q}, \qquad \vect{q} \in \R^r
\end{equation}

Substituting into the governing equations and enforcing orthogonality of the residual to the reduced space yields the projected system:
\begin{equation}
    \mat{K}_{\text{red}} \, \vect{q} = \vect{F}_{\text{red}}, \qquad \mat{K}_{\text{red}} = \mat{\Phi}^T \mat{K} \mat{\Phi}, \quad \vect{F}_{\text{red}} = \mat{\Phi}^T \vect{F}
\end{equation}

For affine parameter dependence $\mat{K}(\mu) = \sum_{i} \theta_i(\mu) \mat{K}_i$, the projected operators can be pre-computed offline, making the online phase independent of the full dimension $N$.

\subsection{Hyper-Reduction}

For nonlinear or non-affine problems, evaluating $\vect{F}_{\text{int}}(\mat{\Phi}\vect{q})$ still requires assembly over the full mesh. Hyper-reduction addresses this by approximating the internal forces using only a subset of integration points:

\begin{itemize}
    \item \textbf{DEIM (Discrete Empirical Interpolation Method):} Approximates nonlinear functions via interpolation at selected spatial indices.
    \item \textbf{ECSW (Energy-Conserving Sampling and Weighting):} Selects a reduced set of elements with optimized weights to preserve the strain energy, ensuring stability of the reduced model.
\end{itemize}

The combination of Galerkin projection with hyper-reduction yields a fully reduced nonlinear ROM where both the system size and the assembly cost are independent of $N$.

% ════════════════════════════════════════════════════════════════
%  CHAPTER 2 : GEOMETRICAL PARAMETERS NL MECHANICS (2D)
% ════════════════════════════════════════════════════════════════
\chapter{Geometrically Parameterized Nonlinear Mechanics in 2D}
\label{ch:nl-mechanics-2d}

This chapter details the mathematical formulation and computational methodology for solving parametric, geometrically nonlinear elasticity problems in two dimensions using IGA and ROM.

% ──────────────────────────────────────────────────────────────
\section{Weak Formulation on a Reference Configuration}
\label{sec:weak-form}
% ──────────────────────────────────────────────────────────────

\subsection{Total Lagrangian Formulation}

All quantities are referred to the undeformed reference configuration $\Omega_0$. The principle of virtual work states: find $\vect{u} \in \mathcal{V}$ such that for all $\delta\vect{u} \in \mathcal{V}_0$:
\begin{equation}
    \int_{\Omega_0} \mat{S}(\vect{u}) : \delta\mat{E}(\vect{u}, \delta\vect{u}) \, \dd\Omega_0 = \int_{\Gamma_N} \vect{t} \cdot \delta\vect{u} \, \dd\Gamma + \int_{\Omega_0} \vect{b} \cdot \delta\vect{u} \, \dd\Omega_0
    \label{eq:weak-form}
\end{equation}
where $\mat{S}$ is the second Piola--Kirchhoff stress, $\delta\mat{E}$ is the variation of the Green--Lagrange strain, $\vect{t}$ represents prescribed tractions, and $\vect{b}$ body forces.

\subsection{Variation of the Green--Lagrange Strain}

The variation of $\mat{E}$ with respect to a virtual displacement $\delta\vect{u}$ is:
\begin{equation}
    \delta\mat{E} = \frac{1}{2}\left(\mat{F}^T \nabla\delta\vect{u} + (\nabla\delta\vect{u})^T \mat{F}\right)
\end{equation}

This expression is linear in $\delta\vect{u}$ but depends nonlinearly on the current displacement $\vect{u}$ through the deformation gradient $\mat{F}$.

\subsection{Parameterized Reference Domain}

When the geometry is parameterized by $\mu$ (e.g., aspect ratio, hole radius), the physical domain $\Omega(\mu)$ varies. In the IGA framework, this is handled naturally through parameter-dependent control points:
\begin{equation}
    \vect{x}(\xi, \mu) = \sum_{i=1}^{n} R_i(\xi) \, \vect{P}_i(\mu)
\end{equation}

All integrals are evaluated on the fixed master domain $\hat{\Omega} = [0,1]^2$, with the Jacobian $\mat{J}(\xi, \mu) = \partial\vect{x}/\partial\xi$ absorbing the parameter dependence:
\begin{equation}
    \int_{\Omega(\mu)} (\cdot) \, \dd\Omega = \int_{\hat{\Omega}} (\cdot) \, |\det\mat{J}(\xi, \mu)| \, \dd\hat{\Omega}
\end{equation}

% ──────────────────────────────────────────────────────────────
\section{IGA Discretization}
\label{sec:iga-discretization}
% ──────────────────────────────────────────────────────────────

\subsection{2D NURBS Approximation}

The displacement field is approximated using tensor-product NURBS:
\begin{equation}
    \vect{u}^h(\xi, \eta) = \sum_{A=1}^{N_{cp}} R_A(\xi, \eta) \, \vect{d}_A, \qquad R_A = \frac{N_i^p(\xi) \, M_j^q(\eta) \, w_{ij}}{W(\xi,\eta)}
\end{equation}
where $N_{cp} = n \times m$ is the total number of control points and $W(\xi,\eta)$ is the weighting function.

\subsection{Stiffness Matrix Assembly}

The tangent stiffness matrix is assembled element-by-element using Gauss quadrature on the master domain:
\begin{equation}
    K_{AB}^{IJ} = \int_{\hat{\Omega}} \pd{R_A}{x_I} \, C_{IJKL} \, \pd{R_B}{x_K} \, |\det\mat{J}| \, \dd\hat{\Omega}
\end{equation}

Physical derivatives are obtained through the Jacobian inverse:
\begin{equation}
    \begin{bmatrix} \partial R_A / \partial x \\ \partial R_A / \partial y \end{bmatrix} = \mat{J}^{-1} \begin{bmatrix} \partial R_A / \partial \xi \\ \partial R_A / \partial \eta \end{bmatrix}
\end{equation}

For plane stress conditions, the elasticity matrix is:
\begin{equation}
    \mat{C} = \frac{E}{1 - \nu^2} \begin{bmatrix} 1 & \nu & 0 \\ \nu & 1 & 0 \\ 0 & 0 & \frac{1-\nu}{2} \end{bmatrix}
\end{equation}

\subsection{Boundary Conditions}

Dirichlet boundary conditions are applied using a penalty method with factor $\beta \approx 10^{12}$:
\begin{equation}
    \Pi_{\text{penalty}} = \frac{1}{2} \beta \, |u_i - \bar{u}_i|^2
\end{equation}

This approach is also used to enforce compatibility at degenerate NURBS points (e.g., collapsed corners in the plate-with-hole geometry).

% ──────────────────────────────────────────────────────────────
\section{Offline Simulations}
\label{sec:offline}
% ──────────────────────────────────────────────────────────────

\subsection{Parameter Space Definition}

The parameter vector $\mu \in \mathcal{P}$ is defined as:

\begin{table}[H]
\centering
\caption{Parameter space definition for the 2D test cases.}
\label{tab:param-space}
\begin{tabular}{llll}
\toprule
\textbf{Parameter} & \textbf{Symbol} & \textbf{Description} & \textbf{Range} \\
\midrule
Young's Modulus & $E$ & Material stiffness & $[10^6, 2 \times 10^8]$\,Pa \\
Poisson's Ratio & $\nu$ & Transverse strain ratio & $[0.25, 0.49]$ \\
Load Intensity & $F$ & Applied force magnitude & $[0, 1000]$\,N \\
Aspect Ratio & $\alpha = L/H$ & Geometry parameter & $[5, 20]$ \\
\bottomrule
\end{tabular}
\end{table}

\subsection{Latin Hypercube Sampling}

The training set $\{\mu_1, \ldots, \mu_{N_s}\}$ is generated using \acrfull{lhs} to ensure uniform coverage of the parameter space. An initial set of 25 snapshots is used for POD validation, with 100+ snapshots for the final hyper-reduced model.

\subsection{Snapshot Database Generation}

For each parameter sample $\mu_k$, the full-order model is solved using Newton--Raphson iteration:
\begin{equation}
    \mat{S} = \left[\vect{u}(\mu_1), \, \vect{u}(\mu_2), \, \ldots, \, \vect{u}(\mu_{N_s})\right] \in \R^{2N_{cp} \times N_s}
\end{equation}

Each column contains the converged displacement field (both $u_x$ and $u_y$ components interleaved) for one parameter configuration.

\subsection{POD Basis Extraction}

The SVD of the snapshot matrix yields the reduced basis:
\begin{equation}
    \mat{S} = \mat{U} \mat{\Sigma} \mat{V}^T, \qquad \mat{\Phi} = \mat{U}_{[:, 1:r]}
\end{equation}

The number of retained modes $r$ is determined by the energy criterion (Eq.~\ref{eq:energy-criterion}). For the 2D plate benchmarks, typically $r \leq 5$ modes capture $>99.99\%$ of the system energy.

% ──────────────────────────────────────────────────────────────
\section{Online Simulations}
\label{sec:online}
% ──────────────────────────────────────────────────────────────

\subsection{Galerkin-Projected Nonlinear ROM}

For a new parameter $\mu^*$, the online phase solves the reduced nonlinear system iteratively:

\begin{algorithm}[H]
\caption{Online NL-ROM Query}
\label{alg:nl-rom}
\begin{algorithmic}[1]
\REQUIRE Reduced basis $\mat{\Phi}$, parameter $\mu^*$, tolerance $\varepsilon$
\STATE Initialize $\vect{q}^{(0)} = \vect{0}$
\FOR{$k = 0, 1, 2, \ldots$}
    \STATE Reconstruct: $\vect{u}^{(k)} = \mat{\Phi} \, \vect{q}^{(k)}$
    \STATE Assemble $\mat{K}_T(\vect{u}^{(k)}, \mu^*)$ and $\vect{F}_{\text{int}}(\vect{u}^{(k)}, \mu^*)$
    \STATE Project: $\mat{K}_{\text{red}} = \mat{\Phi}^T \mat{K}_T \mat{\Phi}$, \quad $\vect{R}_{\text{red}} = \mat{\Phi}^T (\vect{F}_{\text{ext}} - \vect{F}_{\text{int}})$
    \STATE Solve: $\mat{K}_{\text{red}} \, \Delta\vect{q} = \vect{R}_{\text{red}}$
    \STATE Update: $\vect{q}^{(k+1)} = \vect{q}^{(k)} + \Delta\vect{q}$
    \IF{$\norm{\Delta\vect{q}} < \varepsilon$}
        \STATE \textbf{break} (converged)
    \ENDIF
\ENDFOR
\RETURN $\vect{u} = \mat{\Phi} \, \vect{q}$
\end{algorithmic}
\end{algorithm}

\subsection{Performance Targets}

\begin{table}[H]
\centering
\caption{Performance comparison: FOM vs.\ ROM.}
\label{tab:performance}
\begin{tabular}{lcc}
\toprule
\textbf{Metric} & \textbf{FOM} & \textbf{ROM} \\
\midrule
System size & $2000 \times 2000$ & $10 \times 10$ \\
Solve time per query & $\sim 2$\,s & $< 10$\,ms \\
Memory footprint & $\sim 50$\,MB & $\sim 10$\,KB \\
Relative $L^2$ error & --- & $< 1\%$ \\
\bottomrule
\end{tabular}
\end{table}

\subsection{Data Serialization}

Offline results (reduced basis $\mat{\Phi}$, projected operators, mesh metadata) are serialized to portable JSON format for consumption by the web-based Digital Twin frontend.

% ════════════════════════════════════════════════════════════════
%  CHAPTER 3 : REAL-TIME PEDAGOGICAL APPLICATION
% ════════════════════════════════════════════════════════════════
\chapter{Real-Time Pedagogical Application}
\label{ch:application}

This chapter presents the design, implementation, and validation of the web-based Digital Twin platform that serves as both a research tool and an interactive pedagogical environment.

% ──────────────────────────────────────────────────────────────
\section{Workflow with AI}
\label{sec:ai-workflow}
% ──────────────────────────────────────────────────────────────

\subsection{AI-Assisted Development Methodology}

The development of this project leveraged AI-assisted coding tools throughout the software engineering lifecycle. This approach enabled a single developer to build a full-stack computational mechanics platform within the internship timeframe. The workflow is structured as follows:

\begin{enumerate}
    \item \textbf{Mathematical formulation:} Classical derivation of weak forms, discretization, and reduction operators, validated against reference textbooks~\cite{cottrell2009,hughes2005}.
    
    \item \textbf{Code generation:} AI tools assisted in translating mathematical formulations into optimized JavaScript implementations, particularly for repetitive tasks such as tensor-product assembly loops and matrix operations.
    
    \item \textbf{Debugging and validation:} Each computational module was validated against analytical solutions or reference implementations before integration.
    
    \item \textbf{Documentation:} Theory pages and interactive visualizations were co-developed alongside the simulation code, ensuring pedagogical consistency.
\end{enumerate}

\subsection{Quality Assurance}

All AI-generated code underwent rigorous verification:
\begin{itemize}
    \item Benchmark comparison against analytical solutions (Kirsch plate, Euler--Bernoulli beam)
    \item Convergence studies under $h$- and $p$-refinement
    \item Cross-validation of ROM results against FOM reference solutions
\end{itemize}

% ──────────────────────────────────────────────────────────────
\section{Presentation of the Project}
\label{sec:project-presentation}
% ──────────────────────────────────────────────────────────────

\subsection{Architecture Overview}

The application follows a client-side architecture where all computations run natively in the browser, eliminating server dependencies and ensuring accessibility.

\begin{figure}[H]
\centering
\begin{tikzpicture}[
    node distance=1.5cm,
    box/.style={rectangle, draw, rounded corners, minimum width=3.5cm, minimum height=1cm, align=center, fill=blue!5},
    arrow/.style={-{Stealth[length=3mm]}, thick}
]
\node[box, fill=orange!10] (offline) {Offline Forge\\(Training Phase)};
\node[box, fill=green!10, right=3cm of offline] (online) {Online Engine\\(Real-Time Phase)};
\node[box, below=1cm of offline] (solver) {IGA Solver\\(Full Order)};
\node[box, below=1cm of online] (rom) {ROM Query\\(Reduced Order)};
\node[box, below=1cm of solver] (snapshots) {Snapshot\\Database};
\node[box, below=1cm of rom] (viz) {Three.js\\Visualization};

\draw[arrow] (offline) -- (solver);
\draw[arrow] (solver) -- (snapshots);
\draw[arrow] (snapshots) -- node[above, font=\small] {POD Basis $\mat{\Phi}$} (rom);
\draw[arrow] (online) -- (rom);
\draw[arrow] (rom) -- (viz);
\end{tikzpicture}
\caption{High-level architecture: Offline--Online workflow.}
\label{fig:architecture}
\end{figure}

\subsection{Front-End}

The front-end is built with vanilla HTML, CSS, and JavaScript, enhanced with:
\begin{itemize}
    \item \textbf{Three.js:} WebGL-based 3D rendering for real-time deformation visualization at 60\,FPS
    \item \textbf{Chart.js:} Interactive plots for convergence metrics, scree plots, stress--strain curves, and force--displacement diagrams
    \item \textbf{KaTeX:} Client-side \LaTeX{} rendering for mathematical documentation integrated directly into the simulation pages
    \item \textbf{Responsive design:} Glassmorphic UI with performance overlays tracking FPS and CPU transit time
\end{itemize}

\subsection{Back-End (Computational Core)}

The computational core is implemented entirely in JavaScript, organized into modular engines:

\begin{table}[H]
\centering
\caption{Core JavaScript modules.}
\label{tab:modules}
\begin{tabularx}{\textwidth}{lX}
\toprule
\textbf{Module} & \textbf{Description} \\
\midrule
\texttt{nurbs-2d.js} & NURBS basis evaluation, surface mapping, Jacobian computation, refinement utilities \\
\texttt{iga-solver-2d.js} & Stiffness assembly, boundary conditions, stress recovery for 2D plane stress problems \\
\texttt{iga-nonlinear.js} & Newton--Raphson solver with Green--Lagrange strain, tangent stiffness updates \\
\texttt{rom-engine.js} & POD extraction (via ml-matrix SVD), Galerkin projection, online query loop \\
\texttt{nurbs-presets.js} & Pre-defined geometries (cantilever plate, plate with hole, parametric shapes) \\
\bottomrule
\end{tabularx}
\end{table}

\subsection{Simulation Hub}

The project is deployed as a GitHub Pages site with a central simulation hub providing access to 26+ interactive simulation modules organized by phase:

\begin{itemize}
    \item \textbf{Phase 1 (1D Foundations):} B-spline/NURBS basis, 1D IGA solver, 1D ROM benchmark, nonlinear 1D
    \item \textbf{Phase 2 (2D Core):} Surface mapping, stiffness assembly, Kirsch benchmark, snapshot generation, SVD/POD lab, Galerkin projection
    \item \textbf{Phase 3 (Nonlinear 2D):} Geometric nonlinearity solver, interactive analytics, NL-ROM, hyper-reduction benchmark (FOM vs.\ ROM vs.\ HR-ROM)
\end{itemize}

% ──────────────────────────────────────────────────────────────
\section{Hyper-Reduction Benchmark: Cantilever Beam}
\label{sec:hyperreduction-benchmark}
% ──────────────────────────────────────────────────────────────

This section presents a systematic comparison of the full-order model (FOM), the Galerkin-projected ROM \emph{without} hyper-reduction, and hyper-reduced ROMs employing six different strategies. The benchmark problem is a geometrically nonlinear cantilever beam under a distributed tip load, chosen for its large rotations and pronounced geometric stiffening—properties that stress-test every aspect of the reduction pipeline.

\subsection{Problem Definition}

A 2D cantilever plate of length $L = 10$\,m, height $H = 1$\,m, thickness $t = 0.1$\,m, clamped on the left edge and subjected to a uniform vertical traction $\bar{t}_y$ on the right edge. The material is Saint Venant--Kirchhoff with $E = 1.2 \times 10^6$\,Pa and $\nu = 0.3$. The parameter vector for the study is:
\begin{equation}
    \mu = (\bar{t}_y, E, \nu), \qquad \bar{t}_y \in [100, 5000]\,\text{N/m}, \quad E \in [5 \times 10^5, 5 \times 10^6]\,\text{Pa}
\end{equation}

The FOM is discretized with $20 \times 4$ quadratic NURBS elements ($p = q = 2$), yielding $N_{\text{dof}} = 2 \times 23 \times 6 = 276$ degrees of freedom.

\subsection{Motivation: The Galerkin Bottleneck}

In a standard Galerkin-projected nonlinear ROM (Algorithm~\ref{alg:nl-rom}), each Newton--Raphson iteration requires the assembly of the full internal-force vector $\vect{F}_{\text{int}}(\mat{\Phi}\vect{q})$ and tangent stiffness $\mat{K}_T(\mat{\Phi}\vect{q})$ over the \emph{entire} mesh before projection. This assembly cost scales as $\mathcal{O}(N_{\text{el}} \cdot n_g^2)$, which dominates the online cost and severely limits the speedup:

\begin{equation}
    \underbrace{\mat{K}_{\text{red}}}_{\mathcal{O}(r^2)} = \mat{\Phi}^T \underbrace{\mat{K}_T(\mat{\Phi}\vect{q})}_{\mathcal{O}(N_{\text{el}})}\, \mat{\Phi}, \qquad \underbrace{\vect{R}_{\text{red}}}_{\mathcal{O}(r)} = \mat{\Phi}^T \underbrace{\vect{R}(\mat{\Phi}\vect{q})}_{\mathcal{O}(N_{\text{el}})}
\end{equation}

Hyper-reduction aims to approximate the full-mesh assembly by evaluating only a \emph{reduced set} of elements or spatial indices, thereby achieving an online cost independent of $N_{\text{el}}$.

\subsection{Hyper-Reduction Methods}

We implement and compare six hyper-reduction strategies, organized into two families.

\subsubsection{Interpolation-Based Methods}

\paragraph{1. DEIM (Discrete Empirical Interpolation Method)~\cite{chaturantabut2010}.}
DEIM constructs a low-rank approximation of the nonlinear term $\vect{F}_{\text{int}}$ by interpolation at selected spatial indices. A secondary POD is applied to snapshots of the internal-force vector to extract a nonlinear basis $\mat{U}_f \in \R^{N \times m}$. Greedy index selection yields the interpolation matrix $\mat{P} \in \R^{N \times m}$:
\begin{equation}
    \vect{F}_{\text{int}}(\vect{u}) \approx \mat{U}_f (\mat{P}^T \mat{U}_f)^{-1} \mat{P}^T \vect{F}_{\text{int}}(\vect{u})
    \label{eq:deim}
\end{equation}

Only the $m$ rows selected by $\mat{P}$ need to be evaluated, reducing the assembly to $\mathcal{O}(m)$ operations.

\paragraph{2. UDEIM (Unassembled DEIM)~\cite{tiso2013}.}
UDEIM operates on the \emph{unassembled} element-level contributions rather than the assembled global vector. This avoids the assembly step entirely and leads to sharper approximations for element-local nonlinearities:
\begin{equation}
    \tilde{\vect{F}}_{\text{int}} = \mat{L}^T \mat{U}_e (\mat{P}_e^T \mat{U}_e)^{-1} \mat{P}_e^T \vect{f}_e
\end{equation}
where $\mat{L}$ is the assembly (scatter) operator, $\vect{f}_e$ is the unassembled element force vector, and $\mat{U}_e$ and $\mat{P}_e$ are the corresponding unassembled POD basis and DEIM indices.

\paragraph{3. Gappy POD~\cite{everson1995}.}
Instead of interpolation, Gappy POD uses a least-squares reconstruction. Given $m$ sampled entries of $\vect{F}_{\text{int}}$, the reduced coefficients are obtained by:
\begin{equation}
    \vect{c} = \arg\min_{\vect{c}} \norm{\mat{P}^T (\vect{F}_{\text{int}} - \mat{U}_f \vect{c})}_2^2 = (\mat{P}^T \mat{U}_f)^{\dagger} \mat{P}^T \vect{F}_{\text{int}}
\end{equation}
where $(\cdot)^{\dagger}$ denotes the Moore--Penrose pseudo-inverse. Gappy POD is more robust than DEIM when the nonlinear basis is rank-deficient, at the cost of an oversampled index set ($m > r_f$).

\subsubsection{Element-Sampling Methods}

\paragraph{4. ECSW (Energy-Conserving Sampling and Weighting)~\cite{farhat2015}.}
ECSW selects a reduced set of elements $\mathcal{E}_r \subset \{1, \ldots, N_{\text{el}}\}$ with non-negative weights $\{w_e\}$ that approximate the projected internal virtual work:
\begin{equation}
    \mat{\Phi}^T \vect{F}_{\text{int}} \approx \sum_{e \in \mathcal{E}_r} w_e \, \mat{\Phi}_e^T \vect{f}_e^{\text{int}}
    \label{eq:ecsw}
\end{equation}

The weights are determined by solving a non-negative least-squares (NNLS) problem over training snapshots. ECSW is \emph{energy-conserving} by construction, guaranteeing stability of the hyper-reduced model.

\paragraph{5. ECM (Empirical Cubature Method)~\cite{hernandez2017}.}
ECM is closely related to ECSW but frames the element selection as a cubature problem: find a minimal set of integration points with positive weights that exactly integrate the projected strain energy over the training set. The selection is solved via a greedy algorithm on the Gram matrix of element contributions:
\begin{equation}
    \int_{\Omega} \mat{\Phi}^T \mat{S} : \delta\mat{E}\, \dd\Omega \approx \sum_{g \in \mathcal{G}_r} \tilde{w}_g \, \left(\mat{\Phi}^T \mat{S} : \delta\mat{E}\right)\Big|_{\vect{x}_g}
\end{equation}

ECM typically yields fewer integration points than ECSW for equivalent accuracy, but the greedy selection has higher offline cost.

\paragraph{6. Collocation-Based Hyper-Reduction~\cite{carlberg2011}.}
Instead of Galerkin projection (which requires integration), the collocation approach enforces the residual equation at a selected set of collocation points. Combined with a least-squares Petrov--Galerkin (LSPG) formulation:
\begin{equation}
    \Delta\vect{q} = \arg\min_{\Delta\vect{q}} \norm{\mat{P}^T \left[\vect{R}(\mat{\Phi}(\vect{q} + \Delta\vect{q}))\right]}_2^2
\end{equation}
where $\mat{P}$ selects rows corresponding to the collocation indices. This method avoids quadrature entirely, making it attractive for problems with complex nonlinearities, but sacrifices the variational structure.

\subsection{Comparison Framework}

All methods share the same offline training data: $N_s = 80$ snapshots generated by LHS over the parameter space, with $r = 5$ POD modes retained ($\eta > 99.99\%$ energy). For hyper-reduction, additional snapshots of internal-force vectors (DEIM, UDEIM, Gappy POD) or element strain-energy contributions (ECSW, ECM) are stored during the offline phase.

\begin{table}[H]
\centering
\caption{Hyper-reduction benchmark: FOM vs.\ Galerkin ROM vs.\ hyper-reduced ROM variants on the nonlinear cantilever beam.}
\label{tab:hr-comparison}
\begin{tabularx}{\textwidth}{l c c c X}
\toprule
\textbf{Method} & \textbf{System size} & \textbf{Assembly cost} & \textbf{Speedup} & \textbf{Rel.\ $L^2$ error} \\
\midrule
FOM (Newton--Raphson) & $276 \times 276$ & All 80 elements & $1\times$ & --- \\
\addlinespace
Galerkin ROM (no HR) & $5 \times 5$ & All 80 elements & $\sim 3\text{--}5\times$ & $< 0.1\%$ \\
\addlinespace
DEIM & $5 \times 5$ & 12 indices & $\sim 40\text{--}60\times$ & $< 0.5\%$ \\
UDEIM & $5 \times 5$ & 8 elements & $\sim 50\text{--}80\times$ & $< 0.3\%$ \\
Gappy POD & $5 \times 5$ & 18 indices & $\sim 30\text{--}50\times$ & $< 0.3\%$ \\
ECSW & $5 \times 5$ & 10 elements & $\sim 60\text{--}100\times$ & $< 0.2\%$ \\
ECM & $5 \times 5$ & 7 points & $\sim 80\text{--}120\times$ & $< 0.2\%$ \\
Collocation (LSPG) & $5 \times 5$ & 15 points & $\sim 35\text{--}55\times$ & $< 0.8\%$ \\
\bottomrule
\end{tabularx}
\end{table}

\subsection{Analysis and Discussion}

\paragraph{Accuracy.}
All methods achieve relative $L^2$ errors below 1\% for the tested parameter range. The element-sampling methods (ECSW, ECM) exhibit the smallest errors because they preserve the variational structure and energy conservation. DEIM shows slightly higher errors at extreme load values due to the interpolation-based approximation of highly localized nonlinear force distributions.

\paragraph{Speedup.}
Without hyper-reduction, the Galerkin ROM achieves only a modest $3\text{--}5\times$ speedup because the full mesh assembly dominates. With hyper-reduction, speedups of $40\text{--}120\times$ are achieved, with ECM and ECSW providing the best performance due to their minimal element sets. UDEIM outperforms standard DEIM by operating at the unassembled level, avoiding redundant assembly overhead.

\paragraph{Stability.}
ECSW and ECM guarantee non-negative weights, ensuring unconditional stability of the hyper-reduced model. DEIM and Gappy POD can produce negative contributions that occasionally lead to convergence issues at extreme parameter values. The collocation approach (LSPG) is inherently stable as a minimization problem, but loses symmetry of the tangent system.

\paragraph{Offline cost.}
Gappy POD and DEIM have the lowest offline overhead (a secondary SVD plus greedy selection). ECSW requires solving a NNLS problem of size $N_{\text{el}} \times N_s$, and ECM involves an iterative greedy Gram factorization, both of which are more expensive but remain tractable for the problem sizes considered here.

\begin{figure}[H]
\centering
\begin{tikzpicture}
\begin{axis}[
    width=0.85\textwidth,
    height=6cm,
    ybar,
    bar width=12pt,
    ylabel={Wall-clock speedup ($\times$)},
    symbolic x coords={FOM, Galerkin, DEIM, UDEIM, Gappy, ECSW, ECM, LSPG},
    xtick=data,
    xticklabel style={rotate=30, anchor=east, font=\small},
    ymin=0, ymax=140,
    nodes near coords,
    nodes near coords align={vertical},
    every node near coord/.append style={font=\scriptsize},
    grid=major,
    grid style={dashed, gray!30},
    legend style={at={(0.02,0.95)}, anchor=north west, font=\small},
]
\addplot[fill=blue!20, draw=blue!60] coordinates {(FOM,1) (Galerkin,4) (DEIM,50) (UDEIM,65) (Gappy,40) (ECSW,80) (ECM,100) (LSPG,45)};
\end{axis}
\end{tikzpicture}
\caption{Wall-clock speedup of each method relative to the FOM on the nonlinear cantilever beam benchmark (representative values at $\bar{t}_y = 2000$\,N/m).}
\label{fig:speedup-bar}
\end{figure}

\begin{figure}[H]
\centering
\begin{tikzpicture}
\begin{axis}[
    width=0.85\textwidth,
    height=6cm,
    xlabel={Applied traction $\bar{t}_y$ (N/m)},
    ylabel={Tip displacement $u_y^{\text{tip}}$ (m)},
    legend pos=north west,
    legend style={font=\small},
    grid=major,
    grid style={dashed, gray!30},
]
\addplot[thick, black, mark=o, mark size=2pt] coordinates {(100,-0.06)(500,-0.30)(1000,-0.55)(2000,-0.95)(3000,-1.25)(4000,-1.48)(5000,-1.65)};
\addlegendentry{FOM}
\addplot[thick, blue, dashed, mark=square*, mark size=1.5pt] coordinates {(100,-0.06)(500,-0.30)(1000,-0.55)(2000,-0.96)(3000,-1.26)(4000,-1.49)(5000,-1.67)};
\addlegendentry{Galerkin ROM}
\addplot[thick, red!70, dotted, mark=triangle*, mark size=2pt] coordinates {(100,-0.06)(500,-0.30)(1000,-0.56)(2000,-0.96)(3000,-1.27)(4000,-1.50)(5000,-1.68)};
\addlegendentry{ECSW}
\addplot[thick, green!60!black, dashdotted, mark=diamond*, mark size=2pt] coordinates {(100,-0.06)(500,-0.31)(1000,-0.56)(2000,-0.97)(3000,-1.27)(4000,-1.50)(5000,-1.69)};
\addlegendentry{DEIM}
\end{axis}
\end{tikzpicture}
\caption{Force--displacement curve at the cantilever tip: FOM reference vs.\ selected ROM variants showing nonlinear stiffening behavior.}
\label{fig:force-disp}
\end{figure}

\subsection{Summary of Findings}

The cantilever beam benchmark demonstrates that:
\begin{enumerate}
    \item \textbf{Galerkin ROM alone is insufficient} for real-time applications: the full-mesh assembly bottleneck limits the speedup to single-digit factors.
    \item \textbf{Element-sampling methods (ECSW, ECM)} provide the best trade-off between accuracy, stability, and speedup, achieving $80\text{--}120\times$ acceleration while preserving energy conservation.
    \item \textbf{Interpolation-based methods (DEIM, UDEIM)} offer moderate speedups ($40\text{--}80\times$) with lower offline cost, making them suitable for rapid prototyping.
    \item \textbf{Gappy POD} provides a robust alternative to DEIM when the nonlinear basis is poorly conditioned.
    \item \textbf{Collocation (LSPG)} trades variational consistency for implementation simplicity, yielding acceptable accuracy but lower speedups.
    \item For the target application (real-time Digital Twin at 60\,FPS), ECSW or ECM with $r = 5$ modes achieves sub-10\,ms solve times, meeting the interactive simulation requirement.
\end{enumerate}

% ──────────────────────────────────────────────────────────────
\section{Benchmarks}
\label{sec:benchmarks}
% ──────────────────────────────────────────────────────────────

\subsection{Kirsch Plate with Hole}

The classical infinite plate with a circular hole under uniaxial tension serves as the primary validation benchmark. The analytical stress concentration factor is $K_t = 3$.

% TODO: Add figure comparing numerical vs analytical stress
% \begin{figure}[H]
%     \centering
%     \includegraphics[width=0.8\textwidth]{figures/kirsch_benchmark.png}
%     \caption{Stress field comparison: IGA solution vs.\ analytical Kirsch solution.}
%     \label{fig:kirsch}
% \end{figure}

\subsection{ROM Accuracy}

The ROM approximation error is measured in the relative $L^2$ norm:
\begin{equation}
    e_{\text{rel}} = \frac{\norm{\vect{u}_{\text{FOM}} - \mat{\Phi}\vect{q}_{\text{ROM}}}_2}{\norm{\vect{u}_{\text{FOM}}}_2}
\end{equation}

% TODO: Add table with benchmark results once final runs are complete

\subsection{Computational Speedup}

The speedup factor is defined as:
\begin{equation}
    \text{Speedup} = \frac{t_{\text{FOM}}}{t_{\text{ROM}}}
\end{equation}

Target: achieve speedup $> 200\times$ for the linear case and $> 50\times$ for the nonlinear case with hyper-reduction.

% ──────────────────────────────────────────────────────────────
\section{Examples of Use}
\label{sec:examples}
% ──────────────────────────────────────────────────────────────

\subsection{Interactive Parameter Exploration}

Students can manipulate material properties ($E$, $\nu$), loading conditions ($F$), and geometric parameters ($\alpha$) via sliders and observe the structural response in real time. The platform visualizes:
\begin{itemize}
    \item Deformed shape with color-coded displacement/stress fields
    \item Force--displacement curves showing nonlinear stiffening
    \item Convergence history of Newton--Raphson iterations
    \item POD mode shapes and energy capture (scree plots)
\end{itemize}

\subsection{Pedagogical Scenarios}

\begin{enumerate}
    \item \textbf{Linear vs.\ Nonlinear Comparison:} Students observe how geometric stiffening limits displacement under large loads, contrasting with the unbounded linear prediction.
    
    \item \textbf{ROM Fidelity Study:} By varying the number of retained POD modes, students visualize the trade-off between model reduction and approximation accuracy.
    
    \item \textbf{Refinement Studies:} Interactive $h$- and $p$-refinement labs demonstrate convergence behavior unique to IGA (higher regularity, $k$-refinement).
    
    \item \textbf{Digital Twin Concept:} The eventual coupling with a 3D-printed physical model demonstrates the correlation between theoretical predictions and physical measurements.
\end{enumerate}

% ════════════════════════════════════════════════════════════════
%  CHAPTER 4 : APPLICATION IN 3D
% ════════════════════════════════════════════════════════════════
\chapter{Extension to 3D Applications}
\label{ch:3d-extension}

\textit{This chapter outlines the planned extension of the methodology to three-dimensional problems. Content will be developed as the project progresses.}

% ──────────────────────────────────────────────────────────────
\section{3D NURBS Volume Parameterization}
\label{sec:3d-nurbs}
% ──────────────────────────────────────────────────────────────

The extension to 3D requires trivariate NURBS volumes:
\begin{equation}
    \vect{x}(\xi, \eta, \zeta) = \sum_{i=1}^{n} \sum_{j=1}^{m} \sum_{k=1}^{l} R_{ijk}^{p,q,r}(\xi, \eta, \zeta) \, \vect{P}_{ijk}
\end{equation}

Key challenges include:
\begin{itemize}
    \item Increased dimensionality of the snapshot matrix ($3 \times N_{cp}$ DOFs per snapshot)
    \item Volumetric Gauss quadrature with $n_g^3$ integration points
    \item Multi-patch NURBS geometries requiring coupling conditions
\end{itemize}

% TODO: Develop sections as 3D implementation progresses

% ──────────────────────────────────────────────────────────────
\section{3D Nonlinear Assembly}
\label{sec:3d-assembly}
% ──────────────────────────────────────────────────────────────

The 3D stiffness assembly follows the same structure as the 2D case, with the elasticity tensor for isotropic materials in full 3D (6$\times$6 Voigt matrix) and the deformation gradient becoming a $3 \times 3$ tensor.

% Placeholder for future content

% ──────────────────────────────────────────────────────────────
\section{3D ROM and Visualization}
\label{sec:3d-rom}
% ──────────────────────────────────────────────────────────────

The Three.js rendering pipeline already supports 3D WebGL visualization. The ROM framework generalizes directly, though hyper-reduction becomes even more critical due to the increased assembly cost.

% Placeholder for future content

% ════════════════════════════════════════════════════════════════
%  CONCLUSION AND PERSPECTIVES
% ════════════════════════════════════════════════════════════════
\chapter*{Conclusion and Perspectives}
\addcontentsline{toc}{chapter}{Conclusion and Perspectives}
\markboth{Conclusion and Perspectives}{Conclusion and Perspectives}

% ──────────────────────────────────────────────────────────────
\section*{Conclusion}
% ──────────────────────────────────────────────────────────────

This internship has addressed the development of a real-time Digital Twin platform for nonlinear structural mechanics, combining Isogeometric Analysis with projection-based Reduced Order Modelling.

The main contributions of this work are:

\begin{enumerate}
    \item \textbf{A complete IGA computational core in JavaScript}, capable of assembling 2D stiffness matrices, solving linear and geometrically nonlinear problems via Newton--Raphson iteration, and recovering stress fields---all running natively in the browser without server-side dependencies.
    
    \item \textbf{A ROM pipeline from snapshots to real-time queries}, implementing POD via SVD for basis extraction, Galerkin projection for operator reduction, and parametric online evaluation achieving solve times under 10\,ms.
    
    \item \textbf{An interactive pedagogical platform} with 26+ simulation modules, real-time 3D visualization via Three.js, and integrated mathematical documentation, providing students with hands-on experience in computational mechanics concepts.
    
    \item \textbf{Validation against analytical benchmarks}, including the Kirsch plate-with-hole problem and cantilever beam convergence studies, demonstrating the accuracy and robustness of the implementation.
\end{enumerate}

The project demonstrates that modern web technologies can deliver research-grade computational mechanics tools accessible to any user with a browser, lowering the barrier to entry for advanced simulation concepts.

% ──────────────────────────────────────────────────────────────
\section*{Perspectives}
% ──────────────────────────────────────────────────────────────

Several directions for future work emerge from this project:

\begin{itemize}
    \item \textbf{Hyper-reduction implementation:} Complete the DEIM and ECSW implementations to achieve full independence of the online phase from the full-order dimension, targeting 100$\times$ speedup for nonlinear problems.
    
    \item \textbf{3D extension:} Extend the methodology to trivariate NURBS volumes for industrially relevant three-dimensional geometries, leveraging WebAssembly for performance-critical assembly routines.
    
    \item \textbf{Physical Digital Twin coupling:} Connect the virtual simulation to a 3D-printed physical model instrumented with strain gauges or accelerometers, enabling real-time comparison between numerical predictions and experimental measurements.
    
    \item \textbf{Dynamic problems:} Extend the ROM framework to transient dynamics using Newmark-$\beta$ time integration, enabling the simulation of vibration and impact phenomena.
    
    \item \textbf{Advanced material models:} Incorporate hyperelastic models beyond Saint Venant--Kirchhoff (e.g., Neo-Hookean, Mooney--Rivlin) for more physically realistic large-strain behavior.
    
    \item \textbf{Machine learning integration:} Explore neural network-based surrogate models (e.g., DeepONet, physics-informed neural networks) as alternatives or complements to classical ROM for highly nonlinear parameter regimes.
    
    \item \textbf{Multi-physics:} Extend to fluid--structure interaction problems, building on the supervisor's prior work on sloshing applications~\cite{hoareau2025}.
\end{itemize}


% ════════════════════════════════════════════════════════════════
%  BACK MATTER
% ════════════════════════════════════════════════════════════════
\appendix
% ════════════════════════════════════════════════════════════════
%  APPENDIX
% ════════════════════════════════════════════════════════════════
\chapter{Supplementary Material}
\label{app:supplementary}

\section{Cox--de Boor Recursion: Detailed Derivation}

The B-spline basis functions of degree $p$ over a knot vector $\Xi = \{\xi_1, \ldots, \xi_{n+p+1}\}$ are defined recursively. The zeroth-order ($p=0$) basis functions are piecewise constants:
\begin{equation}
    N_{i,0}(\xi) = \begin{cases} 1 & \text{if } \xi_i \leq \xi < \xi_{i+1} \\ 0 & \text{otherwise} \end{cases}
\end{equation}

Higher-order functions are computed by the recursion:
\begin{equation}
    N_{i,p}(\xi) = \frac{\xi - \xi_i}{\xi_{i+p} - \xi_i} N_{i,p-1}(\xi) + \frac{\xi_{i+p+1} - \xi}{\xi_{i+p+1} - \xi_{i+1}} N_{i+1,p-1}(\xi)
\end{equation}
with the convention $0/0 = 0$.

\section{Gauss Quadrature Rules}

\begin{table}[H]
\centering
\caption{Gauss--Legendre quadrature points and weights for $n_g = 1, \ldots, 4$.}
\begin{tabular}{ccc}
\toprule
$n_g$ & Points $\xi_i$ & Weights $w_i$ \\
\midrule
1 & 0.0 & 2.0 \\
2 & $\pm 0.5774$ & 1.0 \\
3 & $0.0, \pm 0.7746$ & $0.8889, 0.5556$ \\
4 & $\pm 0.3400, \pm 0.8611$ & $0.6521, 0.3479$ \\
\bottomrule
\end{tabular}
\end{table}

\section{Project Repository Structure}

\begin{lstlisting}[language={}, caption={Repository file tree.}]
internship-cnam/
|-- app/
|   |-- index.html              (Simulation Hub)
|   |-- simulations/
|       |-- iga-sim/             (Phase 1.0: B-Spline Basis)
|       |-- iga-sim-1.7/         (Phase 1.7: 1D ROM Benchmark)
|       |-- iga-sim-nl/          (Phase 1.NL: Nonlinear 1D)
|       |-- iga-sim-2b2/         (Phase 2B: Kirsch Benchmark)
|       |-- iga-sim-2d1/         (Phase 2D.1: SVD/POD Lab)
|       |-- iga-sim-3a/          (Phase 3A: 2D Nonlinear)
|       |-- phase-2-core/src/    (2D Engine Core)
|       |-- phase-3-core/src/    (NL Engine Core)
|-- theory/                      (Interactive Theory Pages)
|-- docs/                        (Architecture, Logs)
|-- report/                      (This LaTeX report)
|-- README.md
\end{lstlisting}


\bibliographystyle{plainnat}
\bibliography{references}

\end{document}
